\chapter{Inspecter, générer, brancher son éditeur}

Les commandes de ce chapitre ne produisent aucun binaire. Elles répondent à des
questions --- et ce sont elles qu'on tape quand quelque chose surprend.

\section{\texttt{info} --- ce que Jenga a compris}

\begin{nkretenir}
\textbf{C'est la première commande de tout diagnostic}, et son intérêt tient en
une nuance : elle montre ce que Jenga a \emph{compris}, pas ce que vous croyez
avoir écrit.

Quatre lignes de sa sortie tranchent des questions différentes :

\begin{center}
\begin{tabular}{@{}ll@{}}
\toprule
\textbf{Ligne} & \textbf{Elle répond à} \\
\midrule
\texttt{Entry file} & quel workspace a été chargé \\
\texttt{Configurations} & \texttt{Debug} existe-t-il vraiment \\
\texttt{Target OSes} & cette plateforme est-elle déclarée \\
tableau \emph{Projects} & mon projet est-il inclus \\
tableau \emph{Toolchains} & cette chaîne est-elle disponible \emph{ici} \\
\bottomrule
\end{tabular}
\end{center}

La dernière ligne est la seule qui décrive la \textbf{machine} et non la
description. Les autres décrivent votre \texttt{.jenga}.
\end{nkretenir}

\begin{nkpiege}
\textbf{« Projet inconnu » ne se corrige presque jamais dans le fichier du
projet.} La bonne question n'est pas « qu'ai-je mal écrit ? » mais « ce fichier
est-il inclus ? ».

\texttt{jenga info} tranche en une seconde : le projet est dans le tableau, ou il
n'y est pas. On relit sinon un descripteur parfaitement juste, longtemps.
\end{nkpiege}

\section{\texttt{compile-flags} --- les drapeaux réels}

\begin{nkcode}{jenga compile-flags -- capture reelle, extrait}
{"format": "jcdb/1", "compiler": "C:\\msys64\\ucrt64\\bin\\clang++.EXE",
 "msvc": false, "projects": 1}
\end{nkcode}

\begin{nkretenir}
\textbf{C'est le seul outil qui montre ce qu'un filtre a réellement produit.}
Un \texttt{defines} qui n'apparaît pas ici n'a pas été appliqué --- et vous savez
alors que la condition n'a pas mordu, plutôt que de chercher dans le code.

Notez aussi le \emph{chemin} du compilateur. Sur une machine qui en porte
plusieurs, c'est la ligne qui dit lequel travaille --- une question qu'on se pose
rarement et qui explique des écarts inexpliqués.
\end{nkretenir}

\begin{nkretenir}
Le format se destine aux diagnostics d'éditeur : c'est ce que consomme un serveur
de langage pour savoir avec quels drapeaux votre fichier est compilé.

\textbf{Sans lui, l'auto-complétion devine} --- et ses erreurs ne correspondent
pas à celles du compilateur, ce qui est la pire des situations : deux avis, dont
un faux, sans savoir lequel.
\end{nkretenir}

\section{\texttt{ide-setup} --- brancher son éditeur}

\begin{nkcode}{}
jenga ide-setup
jenga ide
\end{nkcode}

\begin{nkretenir}
Elle écrit ce qu'il faut pour que votre éditeur comprenne le projet ---
\texttt{pyrightconfig.json}, le dossier \texttt{.jenga-typings/} du chapitre 2,
et le fichier de réglages \textbf{de chaque éditeur qu'elle détecte}.

\texttt{jenga workspace} l'a déjà fait à la création. On la relance après avoir
ajouté des projets, ou quand l'auto-complétion s'égare.
\end{nkretenir}

\textbf{« Détecte » est le mot qui compte.} Jenga cherche une trace, et n'écrit
que pour ce qu'il trouve :

\begin{center}
\small
\begin{tabular}{@{}ll@{}}
\toprule
\textbf{ce qu'il cherche} & \textbf{ce qu'il en conclut} \\
\midrule
\texttt{.vscode/} & VSCode \\
\texttt{.nkcode/} & NKCode \\
\texttt{.idea/} ou \texttt{*.iml} & JetBrains \\
\texttt{*.sublime-project} & Sublime \\
\texttt{.nvim.lua} ou \texttt{init.lua} & Neovim \\
\midrule
\emph{rien de tout cela} & \texttt{pyrightconfig.json} seul \\
\bottomrule
\end{tabular}
\end{center}

Sans trace, aucun dossier d'éditeur n'est créé --- voir l'encadré du
chapitre~\ref{ch:aaz}. Pour en obtenir un la première fois, on le demande :

\begin{nkcode}{}
jenga ide-setup --editor vscode
jenga ide-setup --editor lsp
jenga ide-setup --info
\end{nkcode}

\begin{nkretenir}
\texttt{-{}-info} \textbf{n'écrit rien}. C'est la forme à employer quand on
veut seulement savoir ce que Jenga croit voir --- et c'est aussi la façon de
vérifier qu'il vous a bien reconnu avant de le laisser écrire.
\end{nkretenir}

\section{\texttt{gen} --- produire un autre système de construction}

\begin{nkcode}{}
jenga gen --cmake
jenga gen --vs
jenga gen --makefile
\end{nkcode}

\begin{nkretenir}
\textbf{C'est un pont, pas le fonctionnement normal.} Il sert quand un outil tiers
--- un analyseur, un profileur, un éditeur --- exige un \texttt{CMakeLists.txt} ou
une solution Visual Studio.

\alerte{} \textbf{Le fichier produit est un instantané.} Il ne se met pas à jour quand
votre \texttt{.jenga} change. Deux descriptions coexistent alors, et elles
divergeront --- c'est mécanique.

\textbf{Régénérez-le, ne l'éditez pas.} Une modification faite dans le fichier
généré sera perdue à la prochaine génération, ou pire : survivra et fera diverger
les deux systèmes sans que personne le sache.
\end{nkretenir}

\begin{nkretenir}
Le filtre \texttt{action:} existe pour ce cas :

\begin{nkcode}{}
with filter("action:gen-cmake"):
    defines(["GENERATION_CMAKE"])
\end{nkcode}

Il permet de décrire quelque chose qui ne vaut que lors d'une génération --- par
exemple contourner une limite de l'outil cible sans polluer la construction
normale.
\end{nkretenir}

\section{Les autres commandes}

\begin{center}
\begin{tabular}{@{}ll@{}}
\toprule
\texttt{install} & installe dépendances et chaînes d'outils locales \\
\texttt{docs} & la documentation \\
\texttt{examples} & les exemples fournis \\
\texttt{config} & la configuration de Jenga \\
\texttt{profile} & profilage \\
\texttt{bench} & mesures \\
\texttt{add} / \texttt{file} & ajoute des fichiers ou des dépendances à un projet \\
\bottomrule
\end{tabular}
\end{center}

\begin{nkretenir}
\textbf{Vingt-sept commandes, et quinze alias d'une lettre} --- \texttt{b} pour
\texttt{build}, \texttt{r} pour \texttt{run}, \texttt{i} pour \texttt{info}\dots

\texttt{jenga help} les énumère toutes. C'est la commande à taper avant de
supposer qu'une chose n'existe pas : \textbf{la liste est plus longue qu'on ne
croit}, et l'on réécrit volontiers à la main ce qu'une commande fait déjà.
\end{nkretenir}

\begin{nkretenir}[Les exemples sont livrés avec Jenga]
\texttt{jenga examples} donne accès aux projets de démonstration ---
\nkfile{Jenga/Exemples/}, une trentaine de dossiers du \emph{hello} console au
projet multi-plateforme.

\textbf{Ils sont plus à jour que n'importe quelle documentation}, parce qu'ils
sont construits par l'intégration continue. Devant un doute sur une fonction,
chercher un exemple qui l'emploie coûte moins qu'une lecture d'API.
\end{nkretenir}

\begin{nkbilan}
Ces commandes ne construisent rien et répondent aux questions les plus fréquentes.
\texttt{info} dit ce que Jenga a compris --- et sa table des chaînes d'outils est
la seule partie qui décrive la machine. \texttt{compile-flags} est le seul outil
qui montre ce qu'un filtre a réellement produit, et le chemin du compilateur qui
travaille. \texttt{gen} produit un instantané qui ne se met pas à jour : on le
régénère, on ne l'édite pas. Et \texttt{jenga help} se tape avant de supposer
qu'une chose n'existe pas.
\end{nkbilan}

\begin{nkquestions}
\begin{enumerate}
\item Quelle partie de la sortie de \texttt{info} décrit la machine et non votre
      description ?
\item Comment savoir qu'un filtre n'a pas mordu ?
\item Pourquoi ne faut-il pas éditer un fichier produit par \texttt{gen} ?
\item À quoi sert le filtre \texttt{action:} ?
\item Combien de commandes Jenga offre-t-il, et comment les lister ?
\end{enumerate}
\end{nkquestions}

\begin{nkpratique}
Sur votre projet, employez \texttt{compile-flags} comme instrument :

\begin{enumerate}
\item relevez le compilateur employé et son chemin complet ;
\item ajoutez un \texttt{defines} sous un filtre \emph{vrai} et vérifiez qu'il
      apparaît ;
\item ajoutez-en un sous un filtre \emph{faux} et vérifiez qu'il n'apparaît pas ;
\item ajoutez-en un sous un filtre \emph{mal orthographié} --- et constatez que
      le résultat est le même qu'au point 3.
\end{enumerate}

Le point 4 est celui qui compte : l'outil ne distingue pas une condition fausse
d'une condition impossible.
\end{nkpratique}

\begin{nkdemo}
Prouvez qu'un fichier généré par \texttt{gen} ne suit pas votre description.

\begin{enumerate}
\item générez un \texttt{CMakeLists.txt} ;
\item ajoutez un projet à votre \texttt{.jenga} ;
\item relancez \texttt{jenga info} --- le nouveau projet apparaît ;
\item relisez le \texttt{CMakeLists.txt} --- il ne l'a pas.
\end{enumerate}

\textbf{Ce que la démonstration établit} : les deux descriptions ont divergé en
une seule modification, et rien ne l'a signalé. C'est ce qui rend un fichier
généré dangereux dès qu'on le garde.
\end{nkdemo}

\begin{nklogique}
Votre éditeur souligne une erreur que le compilateur ne signale pas.

\begin{enumerate}
\item Énumérez trois causes possibles.
\item Laquelle \texttt{jenga compile-flags} permet-elle de confirmer, et
      comment ?
\item Le cas inverse --- le compilateur échoue et l'éditeur ne dit rien --- a-t-il
      les mêmes causes ? Justifiez, puis dites laquelle des deux situations est
      la plus dangereuse et pourquoi.
\end{enumerate}
\end{nklogique}
